% Original: PDF strana 19, gornji slajd.
\slajdid{02-s037}
\begin{frame}[label=02-s037]{Karnoova karta za pet promenljivih}
\setlength{\parskip}{0pt}
% element: 02-s037:e01
\begin{columns}[T,onlytextwidth]
\begin{column}{.47\textwidth}\centering
\Oznaka{Sloj \(E=0\)}\par\vspace{1mm}
\begingroup
\def\DEskala{.68}\def\DEpomak{0}
% Zajednički raspored DC/BA. Parametri DEskala i DEpomak dolaze iz omotača.
% Nakon skaliranja oznake imaju 9 pt, indeksi 7 pt (odobren stil s034).
\pgfmathsetmacro{\DEfont}{9/\DEskala}
\pgfmathsetmacro{\DEindeksfont}{7/\DEskala}
\fontsize{\DEfont}{\DEfont}\selectfont
\renewcommand{\small}{\fontsize{\DEfont}{\DEfont}\selectfont}
\scalebox{\DEskala}{%
\begin{karnaugh-map}(label=middle)[4][4][1][$A$][$B$][$C$][$D$]
\foreach \x/\y/\base in {1/3/0,2/3/1,3/3/3,4/3/2,1/2/4,2/2/5,3/2/7,4/2/6,1/1/12,2/1/13,3/1/15,4/1/14,1/0/8,2/0/9,3/0/11,4/0/10}{
 \pgfmathtruncatemacro{\indeks}{\base+\DEpomak}
 \node[anchor=south east,inner sep=1pt,font=\fontsize{\DEindeksfont}{\DEindeksfont}\selectfont] at (\x,\y) {\indeks};
}
\end{karnaugh-map}}

\endgroup

\end{column}
\begin{column}{.47\textwidth}\centering
\Oznaka{Sloj \(E=1\)}\par\vspace{1mm}
\begingroup
\def\DEskala{.68}\def\DEpomak{16}
% Zajednički raspored DC/BA. Parametri DEskala i DEpomak dolaze iz omotača.
% Nakon skaliranja oznake imaju 9 pt, indeksi 7 pt (odobren stil s034).
\pgfmathsetmacro{\DEfont}{9/\DEskala}
\pgfmathsetmacro{\DEindeksfont}{7/\DEskala}
\fontsize{\DEfont}{\DEfont}\selectfont
\renewcommand{\small}{\fontsize{\DEfont}{\DEfont}\selectfont}
\scalebox{\DEskala}{%
\begin{karnaugh-map}(label=middle)[4][4][1][$A$][$B$][$C$][$D$]
\foreach \x/\y/\base in {1/3/0,2/3/1,3/3/3,4/3/2,1/2/4,2/2/5,3/2/7,4/2/6,1/1/12,2/1/13,3/1/15,4/1/14,1/0/8,2/0/9,3/0/11,4/0/10}{
 \pgfmathtruncatemacro{\indeks}{\base+\DEpomak}
 \node[anchor=south east,inner sep=1pt,font=\fontsize{\DEindeksfont}{\DEindeksfont}\selectfont] at (\x,\y) {\indeks};
}
\end{karnaugh-map}}

\endgroup

\end{column}
\end{columns}
\par\vspace{1mm}
% element: 02-s037:e02
\Oznaka{Susedstvo između slojeva}
Pored susedstva unutar svake karte, u 3D prostoru susedna su i polja koja se poklapaju pri preklapanju ove dve karte: \(0\leftrightarrow16\), \(1\leftrightarrow17\), itd.
\end{frame}
